\documentclass[a4paper,12pt]{article}

% ── Packages ──────────────────────────────────────────────────────────────────
\usepackage[margin=2.5cm]{geometry}
\usepackage{amsmath,amssymb}
\usepackage{graphicx}
\usepackage{natbib}
\usepackage{hyperref}
\usepackage{xcolor}
\usepackage{booktabs}
\usepackage{lineno}
\usepackage{setspace}
\doublespacing
%\linenumbers

% ── Macros ────────────────────────────────────────────────────────────────────
\newcommand{\om}{\omega}
\newcommand{\omunit}{km\,s$^{-1}$\,kpc$^{-1}$}
\newcommand{\vobs}{V_{\rm obs}}
\newcommand{\vadj}{V_{\rm adj}}
\newcommand{\kms}{km\,s$^{-1}$}

% ── Title ─────────────────────────────────────────────────────────────────────
\title{Correction: Cross-Epoch Application of the Omega Kinematic\\
Correction from $z=0$ to $z\sim5$}

\author{
  Flynn, David C.\\
  \small EPS Research, Laurel, MD, USA\\
  \small ORCID: 0000-0002-2768-6650\\
  \small \href{mailto:davidflynn@eps-research.com}{davidflynn@eps-research.com}
}

\date{August 2026}

% ══════════════════════════════════════════════════════════════════════════════
\begin{document}
\maketitle

% ── Correction Notice ────────────────────────────────────────────────────────
\begin{center}
\fbox{\parbox{0.9\textwidth}{
\textbf{Correction Notice --- August 2026.}
This document supersedes Version~1.0 of ``Kinematic Sign Reversal of
the Omega Correction Across Cosmic Time'' (Zenodo record 21383809,
July 2026). An external formula audit identified that the
intermediate- and high-redshift omega values in Version~1.0 were
computed using a mis-parenthesized variant of the canonical equation,
producing spurious negative values. Additionally, the KROSS
intermediate-redshift values were derived from a fixed template
rescaling rather than independently measured rotation-curve boundary
points. The reported cross-epoch sign reversal does not survive
correction. The $z=0$ SPARC result is unaffected. Full details are
given below.
}}
\end{center}

\vspace{1em}

% ── Abstract ──────────────────────────────────────────────────────────────────
\begin{abstract}
We report a correction to our cross-epoch application of the
Flynn \& Cannaliato (2025) kinematic omega parameter $\om$. The
previously reported sign reversal from positive $\om$ at $z=0$ to
uniformly negative $\om$ at $z\gtrsim0.6$ was the result of two
errors: (1)~an operator-precedence parenthesization discrepancy
between the canonical equation and its implementation in the
intermediate- and high-redshift corpus-generation pipelines, and
(2)~the use of template-derived proxy values rather than independently
measured rotation-curve boundary points for the KROSS
intermediate-redshift sample.

Under uniform application of the canonical Equation~6 of
Flynn \& Cannaliato (2025), the $z=0$ SPARC result is unchanged
(median $\om = +7.13$\,\omunit, $N=175$, 100\% positive). The
High-$z$ Corpus Z1 values ($N=8$ ALPINE rotators) require
recomputation from their 3DBarolo boundary rings under the corrected
formula; preliminary recomputation yields all-positive values. The
KROSS intermediate-redshift omega values are withdrawn because the
IntZ corpus contains no independently tabulated rotation-curve
boundary points; the stored omega was a deterministic rescaling of
$V_c$ and $R_{\rm half}$, not an independent kinematic measurement.

The $z=0$ peer-reviewed omega correction framework
\citep{flynnCannaliato2025} and the 84-galaxy baryonic validation
\citep{flynn2026baryonic} are not affected by this correction.
The cross-epoch sign-reversal claim, including all associated
conclusions regarding a disk assembly phase transition, is withdrawn.
\end{abstract}

\textbf{Keywords:} galaxies: kinematics and dynamics --- errata ---
galaxies: evolution --- galaxies: high-redshift

% ── 1. Summary of the Error ──────────────────────────────────────────────────
\section{Summary of the Error}
\label{sec:error}

The canonical omega correction parameter introduced by
Flynn \& Cannaliato (2025) is defined as
\begin{equation}
  \om = \frac{V_2}{R_2} - \frac{V_1}{R_1}\left(\frac{R_1}{R_2}\right)^{3/2},
  \label{eq:omega}
\end{equation}
where $R_1$ and $R_2$ are the innermost and outermost measured radii,
and $V_1$ and $V_2$ are the corresponding observed rotation velocities.
The parameter carries dimensions of angular frequency; throughout this
work it is expressed in \omunit{} (noting that
$1$\,km\,s$^{-1}$\,kpc$^{-1}$ $\approx 1.023$\,rad\,Gyr$^{-1}$).

The $z=0$ SPARC pipeline correctly implements
Equation~\ref{eq:omega}. However, the intermediate-redshift (IntZ)
and high-redshift (Z1) corpus-generation code implemented a
parenthesized variant:
\begin{equation}
  \om_{\rm err} = \left(\frac{V_2}{R_2} - \frac{V_1}{R_1}\right)
                  \left(\frac{R_1}{R_2}\right)^{3/2}.
  \label{eq:omega_wrong}
\end{equation}
These two expressions are not algebraically equivalent.
Equation~\ref{eq:omega} applies the $(R_1/R_2)^{3/2}$ scaling only
to the inner term, whereas Equation~\ref{eq:omega_wrong} applies it
to the entire difference. For galaxies with flat or gently falling
rotation curves, Equation~\ref{eq:omega} produces positive values
while Equation~\ref{eq:omega_wrong} can produce negative values from
the same input data.

An operator-precedence correction was applied to 43 repository files
on 12~July 2026, but the corrected formula was not propagated into
the stored corpus data products (JSON, CSV, JSONL) in the IntZ and
Z1 directories, nor into the derived figures, notebooks, or
manuscript statistics. The Version~1.0 preprint therefore compared
correctly computed $z=0$ values against incorrectly computed
$z>0$ values, producing the spurious sign reversal.

\subsection{The KROSS Template-Derivation Problem}

Independent of the parenthesization error, an audit revealed that the
166 KROSS Tier-1 omega values in the IntZ corpus were not computed
from independently tabulated rotation-curve boundary points. Instead,
each value was derived from a fixed power-law template:
$R_1 = 0.3\,R_{\rm half}$, $V_1 = 0.6577\,V_c$, $R_2 = R_{\rm half}$,
$V_2 = V_c$. This produces the deterministic relationship
$\om \propto V_c / R_{\rm half}$ with a constant of proportionality
that is invariant across all 166 galaxies (standard deviation of the
exponent $\sim 10^{-4}$), and with null uncertainty on every record.

Version~1.0 stated that the IntZ omega values used ``the innermost
and outermost tabulated ring radii.'' This statement was incorrect;
the IntZ corpus ships no per-ring $(R, V)$ rotation-curve data. The
KROSS omega quantity is therefore not an independent kinematic
measurement and cannot participate in a like-for-like cross-epoch
comparison with the SPARC or Z1 boundary-ring values.

% ── 2. Corrected Status ──────────────────────────────────────────────────────
\section{Corrected Status of Each Sample}
\label{sec:status}

\begin{table}[h]
\centering
\caption{Cross-epoch omega status after correction.}
\label{tab:status}
\begin{tabular}{llll}
\toprule
Sample & $N$ & Status & Notes \\
\midrule
SPARC ($z=0$, HI) & 175 &
  Retained; $\om = +7.13^{+3.31}_{-2.22}$\,\omunit &
  Canonical Eq.~\ref{eq:omega}; unaffected \\
KROSS ($z\sim0.9$, H$\alpha$) & 166 &
  Withdrawn &
  Template-derived; no boundary-ring data \\
ALPINE Z1 ($z\sim5$, [CII]) & 8 &
  Recomputation in progress &
  3DBarolo boundary rings available \\
\midrule
Cross-epoch sign reversal & --- &
  Withdrawn &
  Artifact of inconsistent formula \\
\bottomrule
\end{tabular}
\end{table}

\subsection{SPARC ($z=0$)}

The SPARC omega values were computed using the canonical
Equation~\ref{eq:omega} and are unaffected by the correction.
The median $\om = +7.13$\,\omunit{} with IQR
$[+4.91, +10.44]$\,\omunit{} across 175 galaxies (100\% positive)
is confirmed by independent reproduction \citep{flynn2026baryonic}.
We note that the Version~1.0 cross-epoch table reported SPARC
$\om = +7.06$, which is the sample \emph{mean}; the correct
\emph{median} is $+7.13$.

\subsection{KROSS ($z\sim0.9$)}

The previously reported median $\om = -9.09$\,rad\,Gyr$^{-1}$
is withdrawn. The IntZ corpus omega field will be set to null with
status \texttt{not\_computable\_from\_current\_corpus} in the
corrected Zenodo release. The legacy value is preserved in an
explicitly deprecated field for audit purposes.

To restore KROSS to the cross-epoch comparison, independently
observed spatially-resolved rotation-curve boundary points
($R_1, V_1, R_2, V_2$ with uncertainties) would need to be
extracted from the KROSS survey data products. This is beyond the
scope of the current correction.

\subsection{ALPINE Z1 ($z\sim5$)}

The Z1 corpus contains genuine per-ring 3DBarolo rotation curves for
8 confirmed Tier-1 rotators. These provide independently measured
boundary points suitable for Equation~\ref{eq:omega}. The previously
reported median $\om = -13.05$\,rad\,Gyr$^{-1}$ was computed using
the incorrect Equation~\ref{eq:omega_wrong}. A preliminary external
recomputation using the canonical equation yields all-positive values
with an approximate median of $+12.3$\,\omunit. The corrected Z1
corpus release will contain per-galaxy recomputed values with full
provenance fields.

% ── 3. Withdrawn Claims ─────────────────────────────────────────────────────
\section{Withdrawn Claims}
\label{sec:withdrawn}

The following claims from Version~1.0 are withdrawn in their entirety:

\begin{enumerate}
  \item The reported monotonic sign reversal from positive $\om$ at
        $z=0$ to uniformly negative $\om$ at $z\gtrsim0.6$.
  \item The ``three-epoch omega arc'' ($+7.13 \to -9.09 \to -13.05$).
  \item The $>3\sigma$ statistical separation between the SPARC and
        KROSS omega distributions.
  \item The step-function character of the sign change between
        $z=0$ and $z\sim0.6$.
  \item The mass independence of the sign reversal across
        $\log M_\star/M_\odot = 9.0$--$11.5$.
  \item The physical interpretation of negative $\om$ as evidence
        for a disk assembly phase transition.
  \item The prediction that RAMSES simulations initialized at $z=6$
        with negative omega initial conditions should reproduce the
        sign reversal at $z\sim0.5$--$0.6$.
  \item All five deposited figures
        (intz\_nb11, nb14, nb15, nb16, nb18), which were products
        of the invalid high-redshift values.
  \item The total galaxy count of 258 in the conclusions (an artifact
        of mixing the 84-galaxy validation sample with the 175-galaxy
        exploratory sample).
\end{enumerate}

% ── 4. What Remains Valid ────────────────────────────────────────────────────
\section{What Remains Valid}
\label{sec:valid}

The following results are not affected by this correction:

\begin{enumerate}
  \item The canonical omega correction equation
        (Equation~\ref{eq:omega}), as published in the peer-reviewed
        Flynn \& Cannaliato (2025).
  \item The $z=0$ SPARC omega result: median $+7.13$\,\omunit{}
        across 175 galaxies, 100\% positive.
  \item The 84-galaxy baryonic validation demonstrating 93\% agreement
        between the omega correction and baryonic rotation velocities
        \citep{flynn2026baryonic}.
  \item All underlying survey data: SPARC HI rotation curves,
        KROSS/KMOS$^{\rm 3D}$ source parameters ($V_c$, $R_{\rm half}$,
        stellar mass, redshift), and ALPINE per-ring 3DBarolo
        measurements.
  \item The EPS Research Astro-RAG Platform corpora (HI v7.0, Dwarf v1.0,
        GC v1.3.2) except for the omega-derived fields in IntZ and Z1.
\end{enumerate}

% ── 5. Corrective Actions ───────────────────────────────────────────────────
\section{Corrective Actions Taken}
\label{sec:actions}

\begin{enumerate}
  \item This correction note supersedes Version~1.0 at Zenodo
        record 21383809.
  \item The IntZ corpus (Zenodo 20453189) will be updated to set
        KROSS omega to null with explicit withdrawal metadata.
  \item The Z1 corpus (Zenodo 21327061) will be updated with
        recomputed omega values using Equation~\ref{eq:omega}.
  \item The platform repository
        (\url{https://github.com/eps-research/rag-corpus-series})
        has been updated to remove sign-reversal claims from the
        README, QuickStart, and affected notebooks.
  \item A centralized canonical omega function with regression
        tests (including the M33 anchor: $\om = 5.093$\,\omunit)
        has been added to prevent recurrence.
  \item A \texttt{CORRECTIONS.md} file has been added to the
        repository root documenting the full correction history.
\end{enumerate}

% ── 6. Lessons and Path Forward ──────────────────────────────────────────────
\section{Lessons and Path Forward}
\label{sec:forward}

This correction illustrates a failure mode in computational
astrophysics pipelines: a formula can be correctly stated in a
manuscript while being incorrectly implemented in the code that
generates the claimed results. The 43-file operator-precedence
patch applied on 12~July 2026 corrected the formula in code but did
not regenerate the stored data products, creating a state in which
corrected code coexisted with uncorrected outputs. Regression tests
anchored to known physical values (e.g., the M33 result) and
CI checks for disallowed formula patterns are now in place to
prevent recurrence.

A legitimate cross-epoch omega comparison requires: (a)~uniform
application of the canonical equation, (b)~independently measured
rotation-curve boundary points at each epoch, and (c)~explicit
accounting for tracer-dependent radius ratios. Condition (a) is now
satisfied. Condition (b) is satisfied for SPARC and Z1 but not for
KROSS. If spatially resolved boundary data can be extracted from
the KROSS survey products in future work, a corrected cross-epoch
analysis may become possible.

% ── Acknowledgements ─────────────────────────────────────────────────────────
\section*{Acknowledgements}

The author thanks Jim Cannaliato for the foundational work on the
omega correction framework \citep{flynnCannaliato2025}. The formula
discrepancy was identified by an external audit conducted by
Dennis Mungai using the Kimi model. The author is grateful for
this finding, which led to the present correction. This research
was conducted independently at EPS Research, Laurel, MD, with no
external funding.

% ── AI-Use Disclosure ────────────────────────────────────────────────────────
\section*{AI-Use Disclosure}
Portions of this manuscript were prepared with assistance from
large-language-model (LLM) systems (Claude, Anthropic; Kimi, Moonshot AI)
for drafting, editing, code auditing, and formula verification.
The formula discrepancy was identified by the Kimi model during an
external audit. All scientific calculations, statistical analyses,
and corrective actions were performed using human-supervised
Python/JupyterLab workflows.

% ── Data Availability ────────────────────────────────────────────────────────
\section*{Data Availability}
All data used in this analysis are publicly available:
\begin{itemize}
  \item Unified HI Rotation Curve Corpus v7.0:
        \url{https://doi.org/10.5281/zenodo.19563417}
  \item IntZ Kinematic Corpus v1.0 (corrected release pending):
        \url{https://doi.org/10.5281/zenodo.20453189}
  \item High-$z$ Kinematic Corpus Z1 (corrected release pending):
        \url{https://doi.org/10.5281/zenodo.21327061}
  \item Correction history:
        \url{https://github.com/eps-research/rag-corpus-series/blob/main/CORRECTIONS.md}
\end{itemize}

% ── References ───────────────────────────────────────────────────────────────
\bibliographystyle{mnras}
\begin{thebibliography}{99}

\bibitem[Flynn \& Cannaliato(2025)]{flynnCannaliato2025}
Flynn, D.~C., \& Cannaliato, J.\ 2025, Frontiers in Astronomy and Space
Sciences, 12, DOI: 10.3389/fspas.2025.1680387

\bibitem[Flynn(2026a)]{flynn2026hi}
Flynn, D.~C.\ 2026a, Unified Galaxy HI Rotation Curve Corpus v7.0,
Zenodo, DOI: 10.5281/zenodo.19563417; arXiv:2604.13489

\bibitem[Flynn(2026b)]{flynn2026baryonic}
Flynn, D.~C.\ 2026b, 175-Galaxy SPARC Baryonic Validation,
Zenodo, DOI: 10.5281/zenodo.20132805

\bibitem[Flynn(2026c)]{flynn2026intz}
Flynn, D.~C.\ 2026c, EPS Research Intermediate-$z$ Kinematic Corpus v1.0,
Zenodo, DOI: 10.5281/zenodo.20453189

\bibitem[Flynn(2026d)]{flynn2026z1}
Flynn, D.~C.\ 2026d, High-$z$ Kinematic Corpus Z1,
Zenodo, DOI: 10.5281/zenodo.21327061; arXiv:2605.25339

\bibitem[Harrison et al.(2017)]{harrison2017}
Harrison, C.~M., et al.\ 2017, MNRAS, 467, 1965

\bibitem[Jones et al.(2021)]{jones2021}
Jones, G.~C., et al.\ 2021, MNRAS, 507, 3540

\bibitem[Lelli et al.(2016)]{lelli2016}
Lelli, F., McGaugh, S.~S., \& Schombert, J.~M.\ 2016, AJ, 152, 157

\bibitem[Wisnioski et al.(2019)]{wisnioski2019}
Wisnioski, E., et al.\ 2019, ApJ, 886, 124

\end{thebibliography}

\end{document}
